\section{The question hidden inside preservation}
\label{sec:intro}

A mathematical reconstruction of Hegelian sublation makes two
commitments: it selects relations that represent the content, and it
selects tests under which that content counts as preserved. Universal
properties can determine a result once those commitments are fixed.
They do not by themselves establish that the commitments express the
intended mediation. We investigate this distinction through exact
calculations of least resolving levels. The calculations show both how
an expanded test domain can change a least result while preserving every
earlier comparison, and how stability across test domains can leave the
interpretation of the represented relations unsettled.

For the first comparison, the relevant enrichment replaces
sets by homotopy types. A set distinguishes its elements; a homotopy type also
allows identifications between elements, identifications between those
identifications, and further coherent structure. The issue is therefore not
merely whether the same symbols remain available in a more elaborate
foundation. The universal constructions denoted by those symbols can change.
A compatibility condition that is sufficient for sets can leave a detectable
obstruction when paths and higher paths are admitted.

The philosophical importance of the question comes from sublation. In the
\emph{Science of Logic}, Hegel distinguishes the result of mediation from an
empty disappearance. In the remark on the expression \emph{Aufheben}, the
earlier determination persists in a changed status: its immediacy is cancelled,
but it is not simply annihilated \cite{HegelSublation}. Whatever else one takes
this to mean, a mathematical reconstruction should explain which aspect of the
earlier determination remains available and in what sense. The bare word
``preservation'' does not yet provide that explanation.

Lawvere's theory supplies a demanding candidate. A level has both a skeletal
and a sheaf presentation in an ambient topos. A higher level resolves it when
the earlier skeleta become sheaves at the higher level. An Aufhebung is a
\emph{least} resolving level. This is substantially stronger than enclosing two
objects in a third. The construction has a mathematical order, an explicit
factorization condition, and a minimality claim
\cite{KennettEtAl,Menni}. The question we ask is whether such leastness is
unchanged when the coefficients become homotopical.

Our answer is negative in a controlled and quantitative sense. We keep the
finite indexing poset and the initial level fixed. We change only the category
of coefficients, from sets to groupoids and then to successively higher
truncated spaces. Every earlier resolution comparison remains unchanged,
including the construction of the objects being tested. Nevertheless the
least resolving full-subcategory level changes at every stage until the whole
indexing category is required. Agreement on every earlier test does not
therefore secure stability of the least adequate reconstruction. This is the
stronger claim we defend, beyond a general reminder that models have
parameters.

\subsection{The result in advance}

For $N\geq3$, let $\Pcat_N$ be the poset of nonempty proper subsets of
$[N]=\{1,\ldots,N\}$. Place two objects $b<a$ below every element of
$\Pcat_N$, obtaining a poset $\Ccat_N$ with $2^N$ objects. The starting level
is $D=\{b\}$. We work with covariant functor categories, equivalently with
presheaves on the opposite posets. Because $b$ is initial, extension from $D$
produces constant diagrams.

For $k\geq0$, let $\Vk=\Spc_{\leq k}$ be the category of $k$-truncated
spaces. Thus $\mathcal V_0\simeq\Set$, while $\mathcal V_1$ may be presented
by groupoids. A full subposet $R$ containing $b$ resolves $D$ if every constant
$\Vk$-valued diagram on $\Ccat_N$ is a right Kan extension of its restriction
to $R$. The main theorem computes the least such $R$:
\begin{equation}\label{eq:introformula}
\Aup_k(D)=
\begin{cases}
\{b\}\,\cup\,\{S\in\Pcat_N:|S|\geq N-k-2\},&0\leq k\leq N-3,\\
\Ccat_N,&k\geq N-2.
\end{cases}
\end{equation}
For unrestricted spaces the answer is also $\Ccat_N$. The word ``least''
here refers to \emph{all} full subposets containing $b$, not merely to a
preselected chain of rank cuts. The restricted class of levels is explicit;
we do not assume a classification of all higher essential subtoposes.

The proof identifies the comparison map at an object $x$ with
\[
A\longrightarrow\Map\bigl(\B(x\downarrow R),A\bigr).
\]
For a proposed resolver, the omitted information is therefore encoded by a
comma-category nerve. At $a$, the nerve for the first line of
\eqref{eq:introformula} is a wedge of
$\binom{N-1}{k+2}$ spheres of dimension $k+1$. These spheres are invisible
to $k$-types but visible to $(k+1)$-types. The theorem is consequently more
informative than an isolated failure of transfer: it specifies how much
coherence each resolution leaves untested.

\subsection{What is and is not claimed as a contribution}

The topology of simplex skeleta, the formula for a limit of a constant
space-valued diagram, and the relation between truncation and connectivity are
standard. They are proved or recalled below and are not proposed as new
foundational results. The contribution offered here is a worked theory of finite
coefficient-relative resolution profiles: a matching-support formula,
explicit leastness calculations, a linear-size family with a sharp
six-object first jump, a nonconstant extension, and a reverse family
in which changed generation makes the least resolver smaller. The formula is an
elementary consequence of Reedy reconstruction, not a new foundation for
it. The philosophical argument pairs the changed-test calculation with
a reader-based comparison in which coefficient stability fails to settle
the interpretation of relations. These comparisons give distinct,
precise tasks to the requirement of determinate preservation.

The finite calculations have a methodological use. The paper
does not solve the general problem of transporting Aufhebung between arbitrary
topoi. It does not derive the sequence of Hegel's categories, and it does not
identify the sphere in an example with a contradiction in Hegel's text. Its
claim is that the family supplies a counterexample to a tempting methodological
inference: mathematical leastness in one semantic setting does not, by itself,
establish leastness in a richer one.

The relevant background is wider than the vocabulary of Aufhebung.
Truncated cofinality already explains why higher coefficients require more
coherent limit data \cite{Du}. We use that background rather than claiming to
rediscover it. What remains to be assessed by further specialist comparison is
whether the specific finite resolution family has appeared in another form.
The argument for its philosophical significance does not depend on presenting
its elementary ingredients as previously unknown.

The ordinary finite-poset endpoint also meets recent work on minimal initial
functors. Dey and Lesnick describe a least full subposet sufficient for
ordinary limit computation, as well as finer minimality questions allowing
relations to be removed \cite{DeyLesnick}. Our coefficient-indexed profile
should be compared with that work and with truncated cofinality, not
presented as a new general theory of ordinary minimal limit computation.

\section{Hegel, Lawvere, and the burden of a formal reconstruction}
\label{sec:philosophy}

\subsection{Three requirements that should not be collapsed}

The Introduction to the \emph{Science of Logic} connects determinate negation
with a result that has a specific content because of what it negates. Its
advance is not the replacement of one empty designation by another. The
method is also not supposed to be an independently supplied device applied
to an indifferent subject matter \cite{HegelIntro}. These commitments motivate
three distinct questions for a reconstruction.

First, there is a question of \emph{retention}. What survives the transition?
An embedding may preserve objects while changing the morphisms between them.
A quotient may preserve connected components while erasing loops. A logical
translation may preserve provable sentences while adding new inferential
relations. None of these is preservation without qualification. A reconstruction
must say which of these achievements, or which other achievement, is relevant
to the passage under interpretation.

Second, there is a question of \emph{determination}. Why is this result the
result of this earlier content? A construction that can be attached equally
well to any pair of terms may be mathematically legitimate but explanatorily
weak. What matters is not merely that something later can accommodate
something earlier. The earlier structure should contribute to specifying the
later one.

Third, there is a question of \emph{necessity}. What excludes alternatives?
A universal property can answer this question after the ambient category,
admissible constructions, and relevant equivalences have been fixed. Whether
those choices express the philosophical content is an additional question.
Confusing the two questions turns the uniqueness of an answer to a specified
problem into an argument that the problem was specified in the only adequate
way.

These requirements do not amount to a definition of Hegel's method. They
are obligations incurred by a mathematical interpretation that claims to
represent its preservation, determination, or necessity. The distinction is
important. An interpretation can meet one obligation and fail another; it can
also be mathematically productive while making only a modest claim about
Hegel.

\subsection{Why the categorical proposal deserves a serious test}

Lawvere's use of adjoint structures gives the philosophical vocabulary
mathematical consequences. His account of unity and identity of opposites
studies common retractions and their adjoint forms in examples from calculus
and physics \cite{Lawvere1996}. The value of this approach lies in the
possibility of asking whether the proposed relation exists, how it behaves,
and which further structures it entails. The resulting mathematics can resist
an interpretation as well as illustrate it.

For full subcategories of an indexing category, restriction has left and right
Kan extension adjoints. These two extensions produce different presentations
of the same smaller diagram category inside the larger one. A resolving level
requires the presentation supplied on one side of the earlier adjunction to
be accessible through the other side of the later one. Leastness then
constrains how much additional structure is admitted. In this respect,
Lawvere's formulation directly addresses determination: the earlier skeletal
image constrains which levels can count as resolutions.

The formulation also makes a hidden choice visible. A Kan extension is
computed in a category of coefficients. In sets its limits enforce ordinary
compatibility. In spaces the same notation requests coherent compatibility.
The latter includes data that the former cannot detect. The notation can
remain unchanged while the mathematical demand becomes stronger. Our
calculation makes this dependence explicit.

\subsection{A textual benchmark: from becoming to determinate being}

The transition at the beginning of the \emph{Logic} provides a demanding
benchmark. In the account of coming-to-be and ceasing-to-be, each moment's
direction passes into the other. Becoming depends on their distinction,
yet that distinction is itself unstable; its result is not an empty
return to nothing \cite{HegelBecoming}. The subsequent remark makes
retention more specific: the earlier determinations become moments whose
significance is established within their new unity \cite{HegelSublation}.

Two requirements follow for an interpretation of this particular passage.
It must explain how the earlier content supplies a reason for its
transformation, and how the retained determinations acquire their changed
significance. Merely putting copies of two objects into a third would meet
neither requirement. Exact recovery of an earlier object would meet a
retention condition, but would still not explain its changed status or
the need for the transition.

Our model leaves those two requirements open. Its indexing relations
are supplied in advance, and admitting the next coefficient category is
a decision in the construction. The theorem then forces a different least
resolver. It establishes the conditional necessity of a repair, not the
self-generated necessity of Hegel's transition. The test separates a
recoverability achievement from the stronger interpretive claims that
could otherwise be attached to it.

Hegel's Introduction itself warns against importing mathematical method
into philosophy \cite[online paragraph 62]{HegelIntro}. Our use of
mathematics is diagnostic: it examines a proposed formal correlate. It
does not claim that a calculation generates the movement of the
\emph{Logic}. This addresses the warning without pretending to settle
Hegel's wider assessment of mathematics.

\subsection{The nLab project and its stated restrictions}

nLab's \emph{Science of Logic} entry develops a modal and type-theoretic
reconstruction: concepts are treated as types, moments through modalities,
and opposition through adjoint structures. It connects an abstract treatment
with universes and categorical models. It also places restrictions on its own
results, distinguishing a 0-type argument about determinate negation from
stable-type constructions \cite{nLabLogic}. The separate Aufhebung entry
states a least resolving-level condition \cite{nLabAufhebung}.

These restrictions are examples of the discipline advocated here, not
errors that our finite example corrects. Likewise, nLab's claim that
adjoints are determined is expressly conditional on a concrete interpreted
moment. The question is how to extend comparable explicitness to the
coefficient dependence of least resolving levels. Pr\"ahauser's exposition
distinguishes a broader resolving relation from a unique least
resolution \cite{Praehauser}; terminology alone cannot supply minimality.

The present test is therefore internal to a mathematical possibility the
project itself takes seriously: movement between categorical and homotopical
settings. It is not an objection that a reconstruction fails to preserve every
distinction in every conceivable foundation. Rather, when a reconstruction
uses both ordinary and higher versions of its structures, their agreement
must be established where it is invoked. A restricted theorem remains
restricted even if its notation resembles a more general theorem.

There is a constructive way to respond. One can retain the ordinary
interpretation and justify its restriction to sets. One can strengthen the
higher reconstruction and supply the extra coherence it requires. Or one can
treat the sequence of coefficient categories as a family of distinct
interpretations with distinct explanatory purposes. Each response is more
informative than assuming that a familiar universal property carries its
philosophical force unchanged across the passage.

\subsection{Formalization as a source of philosophical questions}

One might ask whether all this is merely a technical qualification on an
analogy. It would be, if the claim under examination were only that adjunctions
resemble some aspects of opposition. The stakes change when the claim is that
a canonical or necessary transition has been exhibited. At that point,
the conditions of canonicity become part of the philosophical argument.

Hegel's discussion of determinate being offers a useful restraint here. He
distinguishes determinations present in the concept under development from
those supplied by the interpreter's comparison \cite{HegelDasein}. A
homotopical reconstruction need not deny that distinction. It should apply it
to itself. A path, a higher identity, or an ambient infinity-topos is initially
a resource of the reconstruction. Its relevance to the interpreted content
requires an argument. The mathematical theorem can sharpen that argument
by showing exactly what depends on the resource, but cannot supply the
historical or philosophical identification on its own.
